\documentclass[10pt,a4paper]{amsart}
\usepackage[margin=19mm]{geometry}
\usepackage{amsmath,amssymb,mathtools}
\usepackage[scr=rsfs,cal=euler]{mathalfa}
\usepackage{xfrac}
\usepackage[unicode,hidelinks,bookmarks=false]{hyperref}
\newcommand{\ie}{i.e.\ }
\newcommand{\cf}{cf.\ }
\newcommand{\resp}{respectively\ }
\newcommand{\Subsection}{\subsection}
\providecommand{\nobreakdash}{\nobreak\hbox{-}}
\setlength{\emergencystretch}{2em}
\multlinegap0pt
\usepackage{mathtools}
\usepackage{amssymb}
\usepackage{subfig}
\usepackage{tikz}
\usepackage{dsfont}
\usepackage{nicefrac}
\newtheorem{proposition}{Proposition}
\newtheorem{theorem}{Theorem}
\newcommand{\R}{\mathds{R}}
\renewcommand{\Re}{\operatorname{Re}}
\newcommand{\Z}{\mathds{Z}}
\title{Real Riemann Surfaces: Smooth and Discrete}
\author{Johanna D\"untsch, Felix G\"unther}
\date{Source: arXiv:2512.25022v1 (CC BY 4.0)}
\begin{document}
\maketitle

\noindent\emph{Source and licence.} Real Riemann Surfaces: Smooth and Discrete. Version arXiv:2512.25022v1 (CC BY 4.0), distributed by arXiv under the Creative Commons Attribution 4.0 International licence, http://creativecommons.org/licenses/by/4.0/, as recorded in the arXiv licence metadata for this exact version and shown on the abstract page https://arxiv.org/abs/2512.25022v1. Full licence notice: \texttt{licenses/arxiv-2512.25022-CC-BY-4.0.txt}.\par\smallskip

\noindent\textit{Editorial scope.} Counted unit: the article's theorem th:completeperiodmatrix (source0.tex lines 1414--1428). The article's complete original proof of it is delivered with it (source0.tex lines 1430--1493, one \texttt{proof} environment of 64 lines). No mathematics is written, repaired or re-derived: the statement and the proof are the article's own lines, in the article's own notation, and no retained line is substituted or re-flowed.  Retained uncounted as the source-local context the proof needs: lines 476--495; lines 728--734; lines 1072--1074; lines 1373--1392.  Nothing was omitted from any retained span, so the package carries no omission record.  No retained line contains a citation, so the wrapper carries no bibliography and no bibliography entry.  No retained line contains a figure environment, an image-inclusion command or drawing code, so no image is delivered and the package declares no asset dependency.  Editorial formatting only, in addition: an AMS \texttt{amsart} preamble that restates the article's own declarations and macros used by the retained lines, namely the environments \texttt{proposition}, \texttt{theorem} on separate counters, together with the standard packages those lines need.  The retained environments print in this document as Proposition 1, Theorem 1 in order of appearance, whereas the article prints the same items with the numbering of its own full document.  The complete native source of the article as distributed in the arXiv e-print for this exact version is bundled unchanged as \texttt{sources/191960/source0.tex}.
\begin{proposition}\label{prop:realhomology}
	Let $(\Sigma,\tau)$ be a real Riemann surface of genus $g$ with $k$ real ovals. Then, there exists a symplectic basis $(a_1, \ldots, a_g, b_1, \ldots, b_g)$ of $H_1(\Sigma,\Z)$ such that
	\begin{align*}
		\tau(a_i) &= a_i \quad \text{for } 1 \leq i \leq g,\\
		\tau(b_i) &= \sum_{j=1}^g h_{ji} a_j - b_i \quad \text{for } 1 \leq i \leq g,
	\end{align*}
	where $H = (h_{ij})$ is a symmetric matrix with entries in $\{0,1\}$. Furthermore:
	\begin{enumerate}
		\item[(i)] If $\Sigma$ is dividing, then $\mathrm{diag}(H) = 0$ and $\mathrm{rank}(H) = g + 1 - k$.
		\item[(ii)] If $\Sigma$ is non-dividing and $k \neq 0$, then $\mathrm{diag}(H) = 1$ and $\mathrm{rank}(H) = g + 1 - k$.
		\item[(iii)] If $\Sigma$ is non-dividing and $k = 0$, then $\mathrm{diag}(H) = 0$ and
		\[
		\mathrm{rank}(H) =
		\begin{cases}
			g & \text{if } g \equiv 0 \pmod 2, \\
			g - 1 & \text{if } g \equiv 1 \pmod 2.
		\end{cases}
		\]
	\end{enumerate}
\end{proposition}

\begin{proposition}\label{prop:matrix_smooth}
	Let $(\Sigma, \tau)$ be a real Riemann surface. With respect to the homology basis constructed in Proposition~\ref{prop:realhomology}, the period matrix decomposes as
	\[
	\Pi_{\Sigma} = \frac{1}{2} H + i T,
	\]
	where $T \in GL(g, \R)$ is a real matrix and $H \in M(g, \Z_2)$ is the binary matrix describing the action of $\tau$ on the homology basis (from Proposition~\ref{prop:realhomology}).
\end{proposition}

\begin{proposition}\label{prop:holomorphic_existence}
	For any set of $2g$ complex numbers $A_k^B, A_k^W$ ($1 \leq k \leq g$), there exists a unique discrete holomorphic differential $\omega$ with these specified black and white $a$-periods.
\end{proposition}

\begin{theorem}\label{th:realhomology_discrete}
	Let $(\Sigma, \Lambda, \tau)$ be a discrete real Riemann surface of genus $g$ with $k$ discrete real ovals. Then, there exists a symplectic basis $(a_1, \ldots, a_g, b_1, \ldots, b_g)$ of $H_1(X,\Z)$ such that the induced action of $\tau$ on homology is given by:
	\begin{align*}
		\tau(a_i) &= a_i \quad \text{for } 1 \leq i \leq g,\\
		\tau(b_i) &= \sum_{j=1}^g h_{ji} a_j - b_i \quad \text{for } 1 \leq i \leq g,
	\end{align*}
	where $H = (h_{ji})$ is a symmetric matrix with entries in $\{0,1\}$. Furthermore:
	\begin{enumerate}
		\item[(i)] If $(\Sigma, \Lambda, \tau)$ is dividing, then $\mathrm{diag}(H) = 0$ and $\mathrm{rank}(H) = g + 1 - k$.
		\item[(ii)] If $(\Sigma, \Lambda, \tau)$ is non-dividing and $k \neq 0$, then $\mathrm{diag}(H) = 1$ and $\mathrm{rank}(H) = g + 1 - k$.
		\item[(iii)] If $(\Sigma, \Lambda, \tau)$ is non-dividing and $k = 0$, then $\mathrm{diag}(H) = 0$ and
		\[
		\mathrm{rank}(H) =
		\begin{cases}
			g & \text{if } g \text{ is even}, \\
			g - 1 & \text{if } g \text{ is odd}.
		\end{cases}
		\]
	\end{enumerate}
\end{theorem}

\begin{theorem} \label{th:completeperiodmatrix}
	Let $(\Sigma, \Lambda, \tau)$ be a discrete real Riemann surface. Let $\tilde{\Pi}$ be its complete discrete period matrix with respect to the homology basis from Theorem~\ref{th:realhomology_discrete}, and let $H$ be the matrix defined therein.
	\begin{enumerate}
		\item[(i)] \textbf{Type 1 (color-preserving):} If $\tau$ preserves the bipartite coloring (i.e., $\rho_{\tau(Q)} = \overline{\rho}_{Q}$), then:
		\[ 2 \Re(\Pi^{B,B}) = H = 2 \Re(\Pi^{W,W}) \quad \text{and} \quad \Re(\Pi^{W,B}) = 0 = \Re(\Pi^{B,W}). \]
		Consequently, the complete period matrix takes the form:
		\[ \tilde{\Pi} = \begin{pmatrix}
			i \Pi_1 & i \Pi_2 + \frac{1}{2} H \\
			i \Pi_2^T + \frac{1}{2} H & i \Pi_3
		\end{pmatrix}, \]
		where $\Pi_1, \Pi_2, \Pi_3 \in \R^{g \times g}$ are real matrices.
		\item[(ii)] \textbf{Type 2 (color-reversing):} If $\tau$ reverses the bipartite coloring (i.e., $\rho_{\tau(Q)} = 1/\overline{\rho}_{Q}$), then:
		\[ \Pi^{B,B} + \overline{\Pi}^{W,W} = H \quad \text{and} \quad \Pi^{W,B} + \overline{\Pi}^{B,W} = 0. \]
	\end{enumerate}
\end{theorem}

\begin{proof}
	The proof proceeds in analogy to the proof of Proposition~\ref{prop:matrix_smooth}, utilizing the duality of the basis. Let $\{\omega_k^B,\omega_k^W \mid 1 \leq k \leq g\}$ be the basis of discrete holomorphic differentials dual to the homology basis given in Theorem~\ref{th:realhomology_discrete}. Let $h_{ij}$ denote the entries of $H$.
	
	(i) \textbf{Type 1 (color-preserving):} Assume that $\tau$ preserves the bipartite coloring. We define the discrete one-form $\eta_j^B \coloneq \overline{\tau^* \omega_j^B}$.
	
	First, we verify that $\eta_j^B$ is discrete holomorphic. Closedness follows immediately from the closedness of $\omega_j^B$ and the linearity of the pullback and conjugation. For the second condition, we must show that $\eta_j^B$ is locally exact with respect to a discrete holomorphic function.
	Let $Q \in F(\Lambda)$ be a quadrilateral. Since $\omega_j^B$ is discrete holomorphic, there exists a discrete holomorphic function $f$ on $\tau(Q)$ such that $\omega_j^B|_{\tau(Q)} = df$. We define the function $g$ on the vertices of $Q$ by $g(v) \coloneqq \overline{f(\tau(v))}$. Then, for any edge $e=(u,v)$ of the medial graph inside $Q$ (connecting edges of $\Lambda$), we have:
	\[
	\eta_j^B(e) = \overline{\omega_j^B(\tau(e))} = \overline{df(\tau(e))} = \overline{f(\tau(v)) - f(\tau(u))} = g(v) - g(u) = dg(e).
	\]
	Thus, $\eta_j^B|_Q = dg$. It remains to show that $g$ satisfies the discrete Cauchy-Riemann equations on $Q$. Let $b_-, w_-, b_+, w_+$ be the vertices of $Q$ in counterclockwise order. Since $\tau$ reverses orientation, the vertices $\tau(b_-), \tau(w_-), \tau(b_+), \tau(w_+)$ of $\tau(Q)$ are ordered clockwise. The discrete holomorphic function $f$ satisfies the Cauchy-Riemann equation on $\tau(Q)$ with respect to the standard counterclockwise ordering. Traversing $\tau(Q)$ counterclockwise effectively swaps the relative positions of the white vertices, implying:
	\[
	f(\tau(w_-)) - f(\tau(w_+)) = i \rho_{\tau(Q)} (f(\tau(b_+)) - f(\tau(b_-))).
	\]
	Rearranging and taking the complex conjugate yields:
	\[
	\overline{f(\tau(w_+)) - f(\tau(w_-))} = \overline{- i \rho_{\tau(Q)} (f(\tau(b_+)) - f(\tau(b_-)))}.
	\]
	Using the Type 1 condition $\overline{\rho_{\tau(Q)}} = \rho_Q$, we obtain:
	\[
	g(w_+) - g(w_-) = i \rho_Q (g(b_+) - g(b_-)).
	\]
	Thus, $g$ is discrete holomorphic on $Q$, proving that $\eta_j^B$ is a discrete holomorphic differential.
	
	Since $\tau(a_i)=a_i$, the periods of $\eta_j^B$ are:
	\begin{align*}
		2\int_{Ba_i} \eta_j^B &= \overline{2\int_{\tau(Ba_i)} \omega_j^B} = \overline{2\int_{Ba_i} \omega_j^B} = \overline{\delta_{ij}} = \delta_{ij}, \\
		2\int_{Wa_i} \eta_j^B &= \overline{2\int_{\tau(Wa_i)} \omega_j^B} = \overline{2\int_{Wa_i} \omega_j^B} = 0.
	\end{align*}
	By the uniqueness of the canonical basis (Proposition~\ref{prop:holomorphic_existence}), we obtain $\eta_j^B = \omega_j^B$.
	Using the relation $\tau(b_i) = \sum h_{ki} a_k - b_i$, we calculate the $b$-periods:
	\begin{align*}
		\Pi_{ij}^{B,B} &= 2 \int_{Bb_i} \omega_j^B = 2 \int_{Bb_i} \overline{\tau^*\omega_j^B} = \overline{2 \int_{\tau(Bb_i)} \omega_j^B} \\
		&= \overline{2 \sum_{k=1}^{g} h_{ki} \int_{Ba_k} \omega_j^B - 2 \int_{Bb_i} \omega_j^B}
		= \sum_{k=1}^g h_{ki} \delta_{kj} - \overline{\Pi}_{ij}^{B,B} = h_{ij} - \overline{\Pi}_{ij}^{B,B}.
	\end{align*}
	For the off-diagonal block $\Pi^{W,B}$, recalling that the white $a$-periods of $\omega_j^B$ vanish, we get:
	\begin{align*}
		\Pi_{ij}^{W,B} &= 2 \int_{Wb_i} \omega_j^B = \overline{2 \int_{\tau(Wb_i)} \omega_j^B}
		= \overline{2 \sum_{k=1}^{g} h_{ki} \underbrace{\int_{Wa_k} \omega_j^B}_{0} - 2 \int_{Wb_i} \omega_j^B}
		= - \overline{\Pi}_{ij}^{W,B}.
	\end{align*}
	Using $h_{ji}=h_{ij}$, we conclude $2 \Re(\Pi^{B,B}) = H$ and $\Re(\Pi^{W,B}) = 0$.
	Analogously, defining $\eta_j^W \coloneqq \overline{\tau^*\omega_j^W}$ leads to $\eta_j^W = \omega_j^W$, yielding $2 \Re(\Pi^{W,W}) = H$ and $\Re(\Pi^{B,W}) = 0$.
	
	(ii) \textbf{Type 2 (color-reversing):} Assume that $\tau$ does not preserve the bipartite coloring. We define $\eta_j^B \coloneqq \overline{\tau^*\omega_j^W}$.
	
	Using the same logic as in (i), we establish that $\eta_j^B$ is discrete holomorphic. Locally, if $\omega_j^W|_{\tau(Q)} = df$, then $\eta_j^B|_Q = dg$ with $g = \overline{f \circ \tau}$. The check for the Cauchy-Riemann equations differs slightly: $\tau$ swaps black and white vertices, so the black diagonal of $Q$ maps to the white diagonal of $\tau(Q)$. Combined with the orientation reversal and the Type 2 condition $\rho_{\tau(Q)} = 1/\overline{\rho}_Q$, the Cauchy-Riemann relation for $f$ transforms into the correct relation for $g$. This transformation effectively swaps the role of the diagonals and inverts the coefficient.
	
	The period calculation yields:
	\begin{align*}
		2\int_{Ba_i} \eta_j^B &= \overline{2\int_{\tau(Ba_i)} \omega_j^W} = \overline{2\int_{Wa_i} \omega_j^W} = \overline{\delta_{ij}} = \delta_{ij},\\
		2\int_{Wa_i} \eta_j^B &= \overline{2\int_{\tau(Wa_i)} \omega_j^W} = \overline{2\int_{Ba_i} \omega_j^W} = 0.
	\end{align*}
	Thus $\eta_j^B = \omega_j^B$, meaning $\omega_j^B = \overline{\tau^* \omega_j^W}$.
	Calculating the $b$-periods:
	\begin{align*}
		\Pi_{ij}^{B,B} &= 2 \int_{Bb_i} \omega_j^B = \overline{2 \int_{\tau(Bb_i)} \omega_j^W}
		= \overline{2 \sum_{k=1}^{g} h_{ki}\int_{Wa_k}\omega_j^W - 2 \int_{Wb_i} \omega_j^W}
		= h_{ij} - \overline{\Pi}_{ij}^{W,W}.
	\end{align*}
	Similarly, $\Pi_{ij}^{W,B} = -\overline{\Pi}_{ij}^{B,W}$.
	This implies $\Pi^{B,B} + \overline{\Pi}^{W,W} = H$ and $\Pi^{W,B} + \overline{\Pi}^{B,W} = 0$.
\end{proof}

\end{document}
